% REGENERATED from phaseB_result.json (Phase-5 audited plain-language rewrite). Do not hand-edit.
\chapter{Main Empirical Analyses}\label{ch:main}

\section{Data, Sample, and Variable Construction}

We combine five data sources. First, we download Capital~IQ earnings-call transcripts for U.S. public firms over 2002--2018, parsed so that we know each speaker's role and whether a turn falls in the scripted presentation or in the back-and-forth question-and-answer (Q\&A) segment. Second, we retrieve quarterly Compustat North America fundamentals, taking the most recent fiscal quarter that ends on or before the call. Third, we pull CRSP daily records and IBES earnings records, which supply the stock-return, market-return, and earnings-surprise controls that the Section~2.3 decomposition requires. Fourth, we collect SDC merger-and-acquisition records---announcement date, completion or withdrawal date, and the cash-versus-stock split of how each deal is paid for---keeping only deals with a transaction value of at least \$1~million. Fifth, we take Execucomp annual CEO records and aggregate them into a monthly tenure panel that identifies which executive holds the CEO role in each quarter; this identity is the input to the Section~2.3 manager fixed effect. These five sources do not all use the same firm identifier. SDC keys each deal by the acquirer's six-digit CUSIP, while our call panel is keyed by the Compustat firm identifier (gvkey), so we join deals to calls through the acquirer's six-digit CUSIP, using the CRSP--Compustat link table that anchors the master-sample manifest.

We build the language measures in-house. We first tokenize each call---that is, we break it into individual words---and then compute the Loughran--McDonald uncertainty-word share, the fraction of words drawn from a fixed list of uncertainty terms, separately for each speaker and each segment. The CEO's uncertainty share in the Q\&A, UncAnsCEO, is the input that the Section~2.3 first stage (the \citet{dwz} decomposition, estimated on our own sample) turns into the residual measure UncResCEO---the part of the CEO's answer uncertainty that is left after its predictable component is removed. This is the construction step only; the equation itself, and the definition and justification of UncResCEO, live in Section~2.3, and we do not re-derive them here.

The remaining key variables are constructed here but defined in full elsewhere, and we point to each home rather than re-derive it. The residual uncertainty measure UncResCEO is defined in Section~2.3. The pre-announcement indicator PreAnnounceQtr equals one in the single quarter just before a firm's first qualifying deal ($e=-1$) and zero otherwise; it is defined in Sections~2.2 and~2.4. The cash ratio, CashRatio, is cash and equivalents divided by total assets (\texttt{cheq}/\texttt{atq}), defined in Section~2.5 and the Appendix. Cash scrutiny, CashScrutiny, is the share of analyst Q\&A turns that touch on cash or liquidity, counted from the locked word list in Appendix~I; because the median call draws none at all (median $0.0000$), we also use a simple indicator, HighCashScrutiny, that flags any call with cash scrutiny present---only about 11\% of calls clear that bar (mean $0.1127$). Both are defined in Section~2.5. Finally, the firm-financial controls---leverage, log assets, Tobin's~Q, return on assets, capital expenditure, a dividend indicator, and cash-flow volatility---are catalogued in the Appendix.

The estimation samples are built in three layers, which is why the number of observations changes from one table to the next. The first layer is the base call panel. We keep every matched Capital~IQ call of a U.S. public firm that is neither a financial nor a utility company over 2002--2018. We drop any CEO with fewer than five calls, because with so few calls we cannot estimate a speaking style and therefore cannot form the residual. We also restrict the panel to firms that Execucomp covers---roughly the S\&P~1500---because the CEO's identity feeds the Section~2.3 manager fixed effect. The second layer is the SDC deal universe. We keep U.S. public-acquirer deals over 2002--2018 whose status is completed, pending, or withdrawn and whose payment composition is known. Deals that are at least 50\% cash form the cash arm; deals that are at least 50\% stock form the stock comparison arm; and each firm's first qualifying deal sets its event clock. The third layer is a set of design-specific trims. The run-up test drops the post-announcement quarters of treated firms. The event studies cap the post-window at four quarters and truncate it at the firm's next announcement. The scrutiny validity and channel tests require at least three analyst Q\&A turns per call. We state the five-call filter here purely as a construction mechanic; we take up its implications for sample selection in Section~2.4.

Table~\ref{tab:summary_stats} reports summary statistics over 2002--2018. Our earnings-call (language) sample comprises 88,205 earnings-call observations across 1,884 firms; the cash-ratio regressions draw on a broader panel spanning 2,232 firms, because the cash ratio does not require a CEO speaking-style residual and so is not limited to the calls that can form one. The count $N$ varies from row to row, because each variable is summarized only on the calls for which it is available (its complete cases). The residual uncertainty measure UncResCEO is centered near zero, with mean $-0.0000$ and a standard deviation of $0.3010$ (44,900 observations, Panel~B), while the cash ratio averages $0.1708$ with a standard deviation of $0.1806$ (Panel~C). Because the residual can be computed only for CEOs with at least five calls at non-financial and non-utility firms, the uncertainty regressions tend to describe a better-covered set of firms, one that skews toward larger firms, than the cash-ratio regressions do. For that reason every table reports its own firm and observation counts, and we are cautious about reading too much into these summary statistics on their own.

\section{Main Analysis 1: The Pre-Announcement Run-Up}

Our first main analysis (MA1) estimates the single-indicator two-way fixed-effects regression of Section~2.4 four times. The regression carries firm fixed effects and calendar year-quarter fixed effects---in plain terms, it compares a firm against itself and against the same calendar quarter at other firms---and clusters standard errors by firm. We run it once for each outcome---the cash ratio (CashRatio) and the residual uncertainty in the CEO's Q\&A answers (UncResCEO)---in each arm: the cash arm, deals that are at least half cash, and the stock comparison arm, deals that are at least half stock. For treated firms we drop the quarters after the announcement, so the design sees only the run-up into the deal. Identification is within-firm: it sets a treated firm's pre-announcement quarter against that same firm's ordinary quarters. This is a correlational design and does not by itself identify a causal effect; the estimating equation is given in Section~2.4, and we do not restate it here.

Table~\ref{tab:empire_building_did} shows that in the cash arm both measures move in the final quarter before the announcement. Column~2 reports that residual Q\&A uncertainty is elevated, with a coefficient of $0.0461^{***}$ (standard error $0.0172$)---the CEO's answers carry more uncertainty in that pre-deal quarter. Column~1 reports that the cash ratio is rising, with a coefficient of $0.0051^{***}$ (standard error $0.0017$); this coefficient should be read as a change in cash, not a level, because the cash equation carries a partial-adjustment lag of $0.7990^{***}$ (standard error $0.0070$)---that is, most of this quarter's cash is already explained by last quarter's, so the coefficient picks up only the increment added on top. One might worry that the uncertainty result is an artifact of the table's one-tailed reporting convention; it is not. The residual coefficient survives a two-tailed test ($p=.0074$), so the load-bearing result does not depend on the reporting convention. Because each firm is compared with itself, the pattern is within-firm; it remains correlational and does not identify a cause.

The two cash-arm columns are not run on the same set of firms. Because the residual can be formed only for CEOs with at least five calls, the uncertainty test in column~2 uses 27,622 firm-quarters from 1,248 firms, while the cash-ratio test in column~1 uses 67,590 firm-quarters from 2,232 firms. The cash-ratio panel therefore spans more firms (2,232) than the 1,884-firm residual-feasible call sample, because the cash ratio is available for firms whose CEOs lack the five calls the residual requires. The uncertainty test is therefore the more demanding one, run on the smaller universe.

How large are these effects? They are material but modest. The uncertainty coefficient of $0.0461^{***}$ is about fifteen percent of a residual standard deviation---roughly 15.3\%, taking $0.0461$ against the all-universe standard deviation of $0.3010$ reported in Table~\ref{tab:summary_stats}---and it appears in a measure from which everything persistent has already been removed, concentrated in a single quarter that makes up only about one percent of the sample (the pre-announcement indicator has a mean of $0.0094$). The cash coefficient of $0.0051^{***}$ adds roughly half a percentage point of assets to the cash stock---about three percent of the mean cash ratio ($0.0051$ against $0.1708$)---in just one quarter. These magnitudes are supportive but not definitive, and they do not establish a cause.

The stock arm is our comparison. It faces the same disclosure constraint, and its model is almost identical to the cash arm's---the stock partial-adjustment lag of $0.8013^{***}$ (standard error $0.0057$) is nearly the same as the cash arm's $0.7990^{***}$. Yet neither outcome shows a comparable positive run-up. In column~6 residual uncertainty carries a coefficient of $-0.0429$ (standard error $0.0307$), which is not statistically significant; in column~5 the cash ratio carries $-0.0036$, also not significant. We read these as flat, noisy nulls, not as movements in any direction. What differs at $e=-1$, the final quarter before the announcement, is therefore the response, not the model. We should be careful here: reading the two arms side by side is not yet a formal test of the difference between them. That formal cash-versus-stock test is supplied in Section~3.4, our third main analysis.

\section{Main Analysis 2: Differential Timing Around the Announcement}

Our second main analysis (MA2) replaces the single pre-announcement indicator with a four-bin event study. The bins trace the deal through time: \textit{PRE2} marks the quarter two before the announcement ($e=-2$), \textit{PRE1} the quarter just before it ($e=-1$), \textit{GAP} the window in which the deal is announced but not yet closed, and \textit{POST} the quarters after completion. The baseline against which these are measured is the firm's own ordinary quarters ($e\le-3$) together with firms that never acquire. We estimate the specification on the matched universe---the set of firm-quarters where both the residual and the cash ratio are defined, 28,102 firm-quarters across 1,320 firms---and we run both outcomes on exactly these same rows. Because we impose no ordering across the bins, every coefficient is read two-tailed. Putting both outcomes on the identical sample guards the differential-timing reading against the worry that the two paths merely reflect different samples. The design is otherwise that of MA1, and it remains correlational: it does not identify a cause.

Table~\ref{tab:empire_drop_matched} first lets us check for a pre-trend---a drift in the outcomes that is already under way before the deal. There is none on either clock. The \textit{PRE2} coefficient is $0.0068$ (standard error $0.0178$) for the residual and $0.0008$ (standard error $0.0024$) for the cash ratio, neither significant. So the \textit{PRE1} estimates that follow are not simply the tail end of a trend that predates the announcement. We treat this as a validity check, not as a proof of identification.

The same table delivers the discriminating contrast: what happens once the deal becomes public, in the gap window between announcement and closing. On the information clock---the path traced by residual uncertainty, in column~1 of Table~\ref{tab:empire_drop_matched}---the residual spikes at \textit{PRE1} to $0.0473^{***}$ (standard error $0.0178$) and then falls to insignificance the instant the deal is announced, with a \textit{GAP} coefficient of $0.0018$ (standard error $0.0187$), not significant. The drop from \textit{PRE1} to \textit{GAP} is itself significant, $0.0455^{**}$ (standard error $0.0225$), and the full swing from the pre-announcement peak to completion is large and significant, $0.0723^{***}$ (standard error $0.0190$). On the transaction clock---the path traced by the cash ratio, in column~2---the picture differs. Cash is elevated at \textit{PRE1}, $0.0061^{**}$ (standard error $0.0024$), but it does \emph{not} fall when the deal is announced---the \textit{PRE1}-to-\textit{GAP} drop is just $0.0006$, not significant---and it declines only once the cash actually leaves the firm at completion, a \textit{GAP}-to-\textit{POST} drop of $0.0210^{***}$ (standard error $0.0042$) and a \textit{POST} level of $-0.0155^{***}$ (standard error $0.0022$). One caution is in order: the \textit{GAP} cash coefficient itself is $0.0055$ (standard error $0.0034$), positive but not significant, so the persistence of cash through the gap rests on the \emph{absence} of a \textit{PRE1}-to-\textit{GAP} decline, not on a significantly elevated gap level. Throughout, firms are compared with themselves, the design is correlational, and it does not identify a cause.

The entry spike is, in economic terms, the same event we saw in MA1: $0.0473^{***}$ here against $0.0461^{***}$ in Table~\ref{tab:empire_building_did}, about fifteen percent of a residual standard deviation (roughly $15.7\%$ on the $0.3010$ basis). What MA2 adds is the round trip. The swing from the pre-announcement peak to the post-completion quarters, $0.0723^{***}$, says that the uncertainty raised in the run-up is fully unwound once the information becomes public. The cash leg, meanwhile, uses the same well-behaved partial-adjustment specification as MA1, with a lagged-cash coefficient of $0.7547^{***}$ (standard error $0.0108$). This remains a correlational reading, not a causal one.

Column~1 of Table~\ref{tab:empire_drop_placebo} corroborates the round trip on a second sample. On its own complete-case cash universe of 29,535 firm-quarters from 1,326 firms, the cash arm of the comparison event study reproduces the same pattern: the residual again spikes at \textit{PRE1} ($0.0486^{***}$, standard error $0.0174$) and again collapses to insignificance at the gap ($0.0058$, standard error $0.0181$, not significant); the \textit{PRE1}-to-\textit{GAP} drop is significant ($0.0428^{*}$, standard error $0.0221$) and the peak-to-post swing is significant ($0.0681^{***}$, standard error $0.0187$). For now we use only the cash column of this comparison table; its cash-versus-stock contrast is reserved for Section~3.4.

Taken together, the claim rests on the contrast between the two paths, not on the size or sign of any single coefficient. Uncertainty tracks the flow of information and reverses the moment the deal is announced; cash tracks the transaction and falls only at completion---a persistence that is partly mechanical, since the cash does not leave the firm until the deal closes. We do not over-read the one coefficient that dips below baseline, the \textit{POST} uncertainty estimate of $-0.0250^{*}$ (standard error $0.0130$). The reading is correlational and supportive, not definitive; it does not identify a cause, and the mechanism behind it remains open.

\section{Main Analysis 3: Cash-Specificity}

We begin descriptively, rerunning the event study side by side on the stock comparison arm, using the stock column of Table~\ref{tab:empire_drop_placebo} that we held back from Section~3.3. Neither arm trends into the event: the PRE2 coefficient is flat in both ($0.0105$ in the cash arm and $-0.0056$ in the stock arm, neither significant). Only the cash arm then moves. It reproduces the run-up (PRE1 $0.0486^{***}$, standard error $0.0174$; peak-to-post drop $0.0681^{***}$, standard error $0.0187$), while the stock arm shows none of it (PRE1 $-0.0404$, standard error $0.0299$, not significant) and, if anything, its PRE1-to-GAP change is a marginal $-0.0756^{*}$ (standard error $0.0405$), a noisy across-announcement change we do not read as a directional run-up. This is a descriptive side-by-side, not yet the test.

A side-by-side display implies a comparison it never actually runs. So our third main analysis (MA3) puts both treatments into one pooled, interacted model --- the model of Section~2.4 --- and tests the cash-minus-stock difference directly, with a single joint test (a Wald linear restriction, $\beta_c-\beta_s>0$). That restriction is what estimates $\theta_{-1}^{\mathrm{cash}}-\theta_{-1}^{\mathrm{stock}}$ in hypothesis~H1a. This difference, and not a comparison of one significant coefficient against one insignificant one, is the formal cash-specificity test. The pooled baseline is the set of firms that make neither a cash nor a stock deal. The design remains correlational and identifies no causal effect, and it speaks to concentration in cash deals rather than strict specificity.

On the effect itself, column~1 of Table~\ref{tab:empire_cashspec} delivers the formal cash-specificity result. The cash arm's pre-announcement uncertainty coefficient is significant ($0.0459^{**}$, standard error $0.0185$), while the stock arm's is insignificant and negative ($-0.0524$, standard error $0.0436$). The formal cash-minus-stock difference, the Wald restriction, is $0.0983^{**}$ (standard error $0.0476$; $z\approx2.07$, $p=.039$), significant at the five-percent level two-tailed. Because the cash arm is positive while the stock arm is an imprecise null near zero, the gap between them is larger than either coefficient on its own --- roughly a third of a residual standard deviation (about 32.7\% of $0.3010$); the gap therefore rides in part on the imprecise negative stock estimate. We emphasize that this formal linear-restriction test, and not a comparison of one significant against one insignificant coefficient, is what we mean by the cash-specificity claim, and we read it as supportive rather than definitive.

On the proposed cause, the picture is weaker. The cash build-up itself is only faintly cash-sided: the cash-minus-stock difference in the cash ratio is insignificant on the matched universe (column~2, $0.0064$, standard error $0.0071$) and reaches only marginal significance on the full panel (column~3, $0.0092^{*}$, standard error $0.0055$). So while the \emph{effect} concentrates in cash deals, the proposed war-chest mechanism --- that firms deliberately build cash ahead of these deals --- is not established by these tests. We leave that mechanism open.

Overall, we read the result as supported but fragile. One formal test at the five-percent level supports H1a's claim that the run-up concentrates in cash deals --- but its significance rides on the imprecise negative stock estimate ($-0.0524$, standard error $0.0436$), and the cause's own cash-sidedness is weaker still (marginal, and only on the full panel). For that reason we keep our wording at \emph{concentrated} rather than \emph{specific}, we do not treat the difference between a significant and an insignificant coefficient as the test, and we assert no cause.

\chapter{Additional Analyses}\label{ch:additional}

\section{Ruling Out Analyst Scrutiny}

The most immediate rival to our reading is analyst scrutiny. The idea is that the quarters just before a deal draw harder analyst questions about cash, and that it is the hard questions, not the deal itself, that produce the hedged answers we measure. This confound is genuine, not a straw man. Column~4 of Table~\ref{tab:empire_building_did} shows that the incidence of cash scrutiny --- whether a call draws heavy cash questioning at all --- does itself rise in the quarter before a cash acquisition: the coefficient on $\mathrm{HighCashScrutiny}$ is $0.0408^{**}$ (standard error $0.0194$). Scrutiny therefore moves together with both the firm's cash position and the event, and could in principle generate the run-up mechanically. The continuous, volume-based measure tells a weaker story --- in column~3 the $\mathrm{CashScrutiny}$ coefficient is $0.0854$ (standard error $0.0749$), not significant --- so the genuine confound is whether cash scrutiny happens at all (its incidence), not how much of a call it fills (its volume).

We rule it out in three steps. First, we confirm that the scrutiny measure is real --- that it actually moves with how much cash a firm holds. Table~\ref{tab:cash_scrutiny_validity} shows that it does. The share of analyst question-and-answer turns spent on cash and liquidity rises strongly with the firm's cash-to-assets ratio: $0.7530^{***}$ (standard error $0.0949$, one-tailed) without controls, and $0.8519^{***}$ (standard error $0.0967$) once the full control set is added. It also rises with a high-cash indicator, a simple flag for cash-rich firms, $0.1754^{***}$ (standard error $0.0237$). The measure clearly registers how cash-rich a firm is, so the nulls that follow are not the artifact of a dead variable that fails to pick anything up.

Second, scrutiny on its own does not move the residual --- our measure of how hedged a CEO's answers are, taken net of that CEO's persistent speaking style ($\mathrm{UncResCEO}$). In Table~\ref{tab:cash_scrutiny_channel}, the direct effect of scrutiny on that residual is zero to four decimal places (Panel~A, $-0.0000$, standard error $0.0013$, across 41,512 calls), and the same null holds when we run a logit on a dichotomized version of the residual --- a simple high-versus-low split (Panel~B, $-0.0003$, standard error $0.0074$) --- as well as on the raw uncertainty measure and on coarser scrutiny proxies. This test is not circular --- it is not merely the residual predicting itself --- because $\mathrm{CashScrutiny}$ counts the \emph{volume} of cash questions, whereas the uncertainty \emph{wording} of those questions ($\mathrm{UncQue}$) is already netted out of the residual when it is built, in the decomposition of Section~2.3.

Third, and most directly, we put scrutiny and the pre-announcement indicator into the same model, estimated on the cash-acquirer pre-announcement quarters. In Table~\ref{tab:reason_gating}, the pre-announcement indicator stays significant --- $0.0413^{***}$ (standard error $0.0177$) with scrutiny entered alongside it, and $0.0439^{**}$ (standard error $0.0189$) once the interaction is added --- while scrutiny's own effect is null ($0.0000$ and $0.0001$ across the two columns) and the scrutiny-by-pre-announcement interaction is insignificant ($-0.0056$, standard error $0.0111$). Put plainly: the reason for the deal raises uncertainty, the questioning does not, and the questioning does not amplify the reason either. The comparison is within firm --- each treated firm measured against its own ordinary quarters --- and correlational.

We keep the register honest about what this does and does not show. What we have is a failure to find an effect, paired with a dissociation inside a single model --- not a powered equivalence test that could formally rule an effect out. The interaction's confidence interval (roughly $[-0.027,\,+0.016]$) does not exclude effects that would still matter against a run-up near $0.04$; about 89\% of calls draw no cash scrutiny at all; and a contrast between a significant and an insignificant coefficient is not, by itself, a test. So our claim is the narrow one: analyst scrutiny does not account for \emph{this} run-up. We do not claim that scrutiny never matters anywhere. The reading is correlational and supportive, not definitive; it identifies no cause, and the mechanism remains open.

\section{Outsider Reactions: The Bid-Ask Spread}

We now ask whether the speech-uncertainty signal moves the post-call information environment --- how large the information gap is between the firm and outside traders. We proxy that environment with the firm's bid-ask spread, the gap between the price to buy and the price to sell, measured over the 25 trading days that begin on the call date. This is a standard measure of how much outside traders fear trading against better-informed insiders, following \citet{bgt2018}, whose post-call window we adopt, and we decompose the spread by the three speech components of Section~2.3. Two priors frame the test. \citet{dwz} found that the residual uncertainty is unrelated to stock-price or trading-volume responses and explains little of the market reaction, which loads instead on the chief executive's persistent clarity --- but they never examined a liquidity outcome. \citet{bgt2018}, by contrast, found that call language relates to information asymmetry with \emph{opposite} signs across the call's two segments: complexity in the scripted presentation is positively associated with information asymmetry (an obfuscation reading), while complexity in the spontaneous response is negatively associated with it (an information reading).

First, the residual is inert. Across all twelve specifications in Table~\ref{tab:h14c_ceo2_decomp}, UncResCEO carries no association with the post-call spread. In the six contemporaneous columns the coefficients are imprecisely estimated and insignificant, ranging from $-0.0594$ to $-0.1021$ (42,625 calls); one quarter ahead they remain insignificant, ranging from $0.0545$ to $0.0984$ (43,049 calls). Because \citet{dwz} show this residual has no detected link to stock price or trading volume, and never test a liquidity outcome, this extends that finding to the bid-ask spread; the absence of an association here is not evidence of a precisely estimated zero, but it is what we find, and it sits comfortably with \citeauthor{bgt2018}'s result that the spontaneous response segment does not widen information asymmetry.

Second, the scripted presentation is the segment outsiders appear to react to. UncPreCEO is positively associated with the contemporaneous spread, and significantly so in four of the six contemporaneous columns ($0.1664^{**}$, standard error $0.0946$, column~1; $0.3496^{***}$, standard error $0.1253$, column~2; $0.2945^{***}$, standard error $0.1229$, column~4; and $0.1644^{*}$, standard error $0.1184$, column~6; one-tailed). These four significant columns include all three firm-fixed-effect specifications (comparisons made within each firm). UncPreCEO is insignificant in the two industry-fixed-effect specifications with extended controls (columns~3 and~5, $0.0814$ and $0.0108$) and in all six one-quarter-ahead columns, and ClarityCEO is null throughout all twelve. The pattern is consistent with \citeauthor{bgt2018}'s positive relation between scripted-presentation language and information asymmetry (the obfuscation component), and our own bid-ask result locates the contemporaneous spread with the scripted presentation rather than the residual. It does not contradict \citeauthor{dwz}'s finding that the residual is unrelated to stock price and volume, because \citet{dwz} never test the presentation and the bid-ask spread is an information-asymmetry outcome, not a price or volume reaction. We note that this association is contemporaneous only.

Taken together, the contemporaneous spread is associated with the scripted presentation but not with the residual; we do not test the between-segment difference directly. This pattern sits within \citeauthor{dwz}'s own scripted-versus-spontaneous framing and matches \citeauthor{bgt2018}'s presentation-versus-response split. The payoff for our main analysis is one of interpretation rather than identification: because outsiders do not react to the residual, the pre-announcement run-up documented in Section~3 is unlikely to be an artifact of outsider reaction, and is more readily read as a private, insider-side signal. We offer this as a supportive reading, not as proof.

\section{Robustness: Withdrawal as a Resolution Event}

One feature of the Section~3.3 event study deserves a closer look. The post bin is dated to completion: a deal's quarters enter the post bin only once it has closed, and the quarters following a withdrawal are dropped from the event clock. Two of the three round-trip contrasts, peak-to-completion and gap-to-completion, are therefore estimated only over deals that go on to complete. Because completing and withdrawn deals need not be alike, part of the post-completion decline could in principle reflect \emph{which} deals complete rather than the resolution of the disclosure state itself. This is a correlational concern, and we do not claim to identify it away. The peak-to-gap contrast, however, is exempt. It is measured at the public announcement --- a point every announced deal reaches whatever its eventual outcome --- so it conditions on nothing that follows.

The check itself is straightforward. A withdrawal ends the disclosure state as surely as a completion does: once a pending deal is called off and disclosed, the executive is no longer answering around it. So we redefine the post bin to let the quarters at and after a withdrawal enter it alongside the completed quarters, and we re-estimate the Section~2.4 event study on this resolution-inclusive sample, reported in Table~\ref{tab:empire_drop_resolution}. Folding in the withdrawn deals asks a simple question: is the post-resolution decline particular to the deals that complete, or does it appear whenever the disclosure state resolves? We keep the framing modest. We do not claim that a weaker decline here would prove the main result an artifact; the power of this check is limited, as we discuss below. The reading is correlational and supportive rather than definitive.

Table~\ref{tab:empire_drop_resolution} shows that the round trip survives. On the resolution-inclusive sample the residual still falls from its pre-announcement peak to the resolved quarters by $0.0687$ ($p<.01$), against $0.0723$ ($p<.01$) in the main specification. The announcement-quarter leg, from peak to gap, is $0.0457$ ($p<.05$), almost exactly the $0.0455$ ($p<.05$) of the main result. And the cash leg declines from its pre-announcement level to resolution by $0.0204$ ($p<.01$). The sample here covers 28,191 firm-quarters over 1,321 firms. Treating a withdrawn deal as a resolution leaves the round trip in place.

We bound the check honestly. Few first deals are withdrawn and then observed for complete-case quarters inside the post window, so the resolution-inclusive sample adds only 89 firm-quarters and one firm to the 28,102 firm-quarters over 1,320 firms of the main analysis. With so little new data, the post-completion decline can barely move, so what this shows is that the handful of withdrawn deals behave like the completed ones --- not that a completion-driven reading is ruled out on its own. The firmer reassurance is structural: the peak-to-gap leg is measured at announcement and conditions on nothing that follows, so the resolution of the residual at announcement never depended on which deals complete. We therefore read the resolution-inclusive estimates as consistent with the main specification, not as an independent test.

\section{Robustness: The Cash Result Without the Dynamic Term}

The cash-ratio regressions carry the cash ratio's own one-quarter lag alongside firm fixed effects --- a partial-adjustment form, well suited to a sticky stock like cash that moves toward its target only gradually. A dynamic fixed-effects panel of this kind is subject to the Nickell bias, a small-sample bias that the within-firm transformation induces. That bias is not confined to the lag coefficient; it also reaches regressors correlated with lagged cash, so simply declining to interpret the lag does not immunize the event-time bins. The clean guard is therefore to remove the dynamic term altogether, rather than to reason about the size of the bias. The residual takes no lag and is untouched by this, so it serves as our anchor. This is a specification concern, not a claim of causal identification.

So we re-estimate the cash column with static fixed effects, dropping the lagged dependent variable (Table~\ref{tab:empire_drop_staticfe}). Removing the lag changes what the coefficients describe. With the lag, each coefficient is a \emph{change} in the cash ratio around its own recent level --- a partial adjustment. Without it, the coefficients describe the \emph{level} of the cash ratio relative to the baseline. The two specifications are on different scales, so the comparison that carries over is the timing pattern --- the sign and significance of the bins across event time --- not the size of any single coefficient. This is a robustness check and is supportive rather than definitive.

Table~\ref{tab:empire_drop_staticfe} shows that the timing pattern carries over. Without the lag, the cash ratio still does not fall when the deal is announced: the pre-announcement-to-gap change is $0.0012$ below zero and insignificant, much like the $0.0006$ of the with-lag specification. It falls only once the cash leaves the firm at completion, with a gap-to-completion drop of $0.0318$ ($p<.01$) and a completion coefficient of $0.0189$ below zero ($p<.01$). The pre-announcement-to-completion decline is $0.0305$ ($p<.01$), larger than the $0.0216$ ($p<.01$) of the partial-adjustment specification, exactly as expected when the coefficients describe a level rather than a change. The residual column is unchanged, since it never carried the lag, and is identical to Table~\ref{tab:empire_drop_matched}.

How should we read the two side by side? The lag belongs in the main specification, because cash is sticky and partial adjustment is the right description of how it moves; the static specification is not a better model but a check that removes the dynamic term entirely. And because the static specification carries no lagged dependent variable, it carries no dynamic-panel bias --- so the cash timing conclusion (no fall at announcement, a fall at completion) holds whether or not that term is present. The timing of the cash path, on which the differential-timing result rests, does not ride on the dynamic specification. The reading is correlational and supportive rather than definitive.

\section{Robustness: The Main Findings Without the First-Deal Restriction}

Our last robustness check drops the first-deal restriction and stacks every deal a firm makes. The pre-announcement run-up survives. On the all-deals-stacked panel of Table~\ref{tab:empire_building_did}, the cash-arm residual-uncertainty coefficient is $0.0391^{***}$ (standard error $0.0140$) and the cash-ratio coefficient is $0.0033^{**}$ (standard error $0.0015$), while the stock-arm residual coefficient is $-0.0348$ (standard error $0.0272$), a flat null. It also reappears in a binary forward form. Among all firm-quarters, higher CEO Q\&A residual uncertainty is associated with a deal being announced the next quarter, of any payment type: a linear-probability coefficient of $0.0086^{***}$ (standard error $0.0026$, $p=.0011$) and a logit coefficient of $0.3233^{***}$ (standard error $0.0967$, $p=.0008$) in this forward specification (Logit~A), estimated on 40,004 firm-quarters from 1,422 firms with a deal rate of 2.84\% (pseudo-$R^2$ $0.026$). It survives firm and year-quarter fixed effects, with a linear-probability coefficient of $0.0078^{***}$ (standard error $0.00275$; 39,557 firm-quarters, 1,422 firms; within-$R^2$ $0.003$).

The differential-timing round trip survives all-deals as well. On the matched universe, Table~\ref{tab:empire_drop_matched} shows residual uncertainty spiking at PRE1 ($0.0352^{**}$, standard error $0.0148$) and then unwinding: the PRE1-to-GAP drop is $0.0363^{**}$ (standard error $0.0156$) and the PRE1-to-POST drop is $0.0544^{***}$ (standard error $0.0139$). Split by payment type, Table~\ref{tab:empire_drop_placebo} shows the round trip holding in the cash arm (PRE1 $0.0401^{***}$, standard error $0.0145$; PRE1-to-POST drop $0.0543^{***}$, standard error $0.0137$) and not in the stock arm (PRE1 $-0.0272$, standard error $0.0266$, not significant), where the stock arm's PRE1-to-GAP change is $-0.0585^{*}$ (standard error $0.0355$) --- marginally significant and negative, which we do not over-read given the noisy stock null.

The formal cash-concentration result also survives, and in fact strengthens. On the all-deals-stacked panel of Table~\ref{tab:empire_cashspec}, the pooled cash-minus-stock Wald difference is $0.1056^{**}$ (standard error $0.0423$; $z\approx2.50$, $p\approx.013$) --- the cash arm at $0.0447^{**}$ (standard error $0.0180$) against the stock arm's insignificant $-0.0609$ (standard error $0.0385$) --- larger than the first-deal difference of $0.0983^{**}$ ($p=.039$) in the same table. The proposed cause remains the weak leg: the matched cash-ratio cash-minus-stock difference is $0.0071$ (standard error $0.0062$), not significant. So the run-up still concentrates in cash deals --- concentration, not strict specificity, and with no cause asserted --- while the build-up of cash itself is not differentially cash-sided here. In binary form, among deals at $e=-1$, higher residual uncertainty is associated with the deal being a cash deal (coded one) rather than a stock deal (coded zero): a linear-probability coefficient of $0.0613^{**}$ (standard error $0.0283$, $p=.030$) and a logit coefficient of $0.7478^{**}$ (standard error $0.3399$, $p=.028$) in this specification (Logit~B), on 1,105 deal observations from 563 firms (982 cash, 123 stock; cash base rate 88.9\%; pseudo-$R^2$ $0.078$). Under firm and year-quarter fixed effects the coefficient keeps its sign but is no longer significant: a linear-probability coefficient of $0.0644$ (standard error $0.05076$, not significant; $p=.205$; 1,063 deal observations, 563 firms; within-$R^2$ $0.059$). The design stays correlational and identifies no causal effect, and we read it as concentration in cash deals, not strict specificity.
